% 需要在导言区加载：\usepackage{tabularray}

\nopagebreak[4]
\enlargethispage{1cm}
\footnotesize
\begin{longtblr}[
caption = {AI智能体主流攻击工具汇总},
label   = {tab:ai_agent_attack_tools_detail},
halign  = c,
]
{
width   = 0.96\linewidth,
colsep  = 3pt,
colspec = {|>{\rule{0pt}{2pt}}m{1.1cm}
|>{\rule{0pt}{2pt}}m{2.0cm}
|>{\rule{0pt}{2pt}}m{1.2cm}
|>{\rule{0pt}{2pt}}m{\dimexpr0.6\linewidth}
|>{\rule{0pt}{2pt}}X|},
}
\hline
\textbf{工具分类} & \textbf{工具名称} & \textbf{发布日期} &\textbf{核心功能} & \textbf{应用场景} \\
\hline
\SetCell[r=5]{c} 学术\newline 研究型
& AgentLoop \allowbreak ~\cite{Xu2026LoopTrap}
& 2026‑03
& \begin{itemize}[tabitem]
    \item 用于测试智能体在资源消耗、执行失控和拒绝服务场景下的风险；
    \item 诱导LLM智能体错误判断任务未完成，进入循环执行状态。
\end{itemize}
& Agent终止条件投毒与资源放大攻击研究 \\
\cline{2-5}
& PoisonRAG \allowbreak ~\cite{Nazary2025PoisonRAG}
& 2025‑04
& \begin{itemize}[tabitem]
    \item 用于测试RAG系统在恶意数据注入和推荐结果偏移下的安全性；
    \item 对RAG推荐系统的检索数据进行投毒，操控系统返回结果；
    \item 通过篡改项目标签、描述等元数据影响检索增强生成流程。
\end{itemize}
& RAG推荐系统数据投毒与推荐结果操纵研究 \\
\cline{2-5}
& AgentFuzz \allowbreak~\cite{Liu2025AgentFuzz}
& 2025‑03
& \begin{itemize}[tabitem]
    \item 自动检测LLM Agent工具调用链中的污染型漏洞；
    \item 生成面向目标智能体的测试提示与工具调用输入；
    \item 根据执行反馈自适应变异测试样本，提高漏洞触发与复现效率。
\end{itemize}
& 漏洞挖掘、安全测试 \\
% \cline{2-5}
% & MRJ‑Agent ~\cite{Wang2024MRJAgent}
% & 2024‑08
% & \begin{itemize}[tabitem]
%     \item 自动构造多轮对话越狱攻击，测试模型在连续交互中的安全边界；
%     \item 将高风险目标拆解为多轮渐进式请求，诱导模型逐步偏离安全约束；
%     \item 用于评估模型在多轮上下文累积场景下的越狱风险。
% \end{itemize}
% & 安全边界测试、对抗样本研究 \\
% \cline{2-5}
% & PromptInjector~\cite{Liu2024PromptInjection}
% & 2024‑07
% & \begin{itemize}[tabitem]
%     \item 用于对比不同提示注入攻击方式和防御机制的有效性
%     \item 构造提示注入攻击，测试 LLM 应用是否会被外部指令劫持
% \end{itemize}
% & 提示注入攻防实验 \\
\hline
\SetCell[r=7]{c} 开源全能\newline 渗透框架
& MCP‑Strike \allowbreak~\cite{MCPStrike}
& 2026‑06
& \begin{itemize}[tabitem]
    \item 对运行中的MCP Server进行主动式对抗安全测试；
    \item 测试工具描述提示注入、过度敏感参数、路径遍历、响应注入和数据外传诱导等风险；
    \item 借助LLM judge或自适应 LLM智能体对发现结果进行确认、评分和多轮探测。
\end{itemize}
& MCP服务器主动对抗测试与安全评估 \\
\cline{2-5}
& RedAmon \allowbreak~\cite{RedAmon}
& 2026‑06
& \begin{itemize}[tabitem]
    \item 执行侦察、漏洞验证、后渗透分析、AI分诊、代码修复和GitHub PR提的一体化安全测试流程；
    \item 并行运行多个侦察工具，收集子域名、端口、端点等攻击面信息，并写入Neo4j知识图谱。
\end{itemize}
& 授权自动化渗透测试、攻击面分析与漏洞修复闭环 \\
\cline{2-5}
& Strix~\cite{Strix}
& 2025‑10
& \begin{itemize}[tabitem]
    \item 动态运行目标应用代码，发现漏洞并生成PoC对漏洞进行验证；
    \item 编排多个专用Agent并行测试复杂目标，并共享安全发现；
    \item 进行应用安全测试、快速渗透测试和Bug Bounty自动化研究。
\end{itemize}
& 应用安全测试、快速渗透测试、漏洞赏金与CI/CD安全测试 \\
\cline{2-5}
& HexStrike \allowbreak~\cite{HexStrike}
& 2025‑09
& \begin{itemize}[tabitem]
    \item 指挥AI智能体调用150+ 安全工具，执行自动化渗透测试、漏洞发现和安全研究任务；
    \item 编排多个安全智能体协同完成侦察、扫描、利用验证和结果分析；
    \item 对网络、Web 应用、云环境、容器、二进制程序和OSINT目标开展多类型安全测试。
\end{itemize}
& 红队演练、渗透测试 \\
\cline{2-5}
& Ankou~\cite{Ankou}
& 2025‑08
& \begin{itemize}[tabitem]
    \item 提供面向红队演练的模块化C2架构、监听器和agent管理；
    \item 支持多种传输与自定义agent/handler扩展；
    \item 通过AI Operations面板辅助总结结果、发现异常和建议后续；操作。
\end{itemize}
& 授权红队C2、Agent管理与后渗透操作 \\
\cline{2-5}
& RAPTOR \allowbreak~\cite{RAPTOR}
& 2025‑06
& \begin{itemize}[tabitem]
    \item 对代码库或二进制目标进行攻击面分析、数据流追踪和漏洞变体搜索；
    \item 进行 Semgrep/CodeQL 扫描、软件组成分析、模糊测试和崩溃根源分析。
\end{itemize}
& 源代码与二进制漏洞分析、验证及修复 \\
% \cline{2-5}
% & AutoPentest‑DRL~\cite{AutoPentestDRL}
% & 2024‑11
% & \begin{itemize}[tabitem]
%     \item 结合 Nmap 和 Metasploit 对真实网络进行扫描、漏洞验证和路径执行
%     \item 基于深度强化学习为目标网络规划自动化渗透攻击路径
% \end{itemize}
% & 复杂网络环境渗透测试 \\
\hline
\SetCell[r=5]{c} 开源专项\newline 攻击工具
& Deadend CLI~\cite{DeadendCLI}
& 2025‑12
& \begin{itemize}[tabitem]
    \item 提供自托管的Agentic pentest CLI工作流；
    \item 支持通过LiteLLM 接入多种模型，用于黑盒渗透测试任务；
    \item 在XBOW validation suite等基准上评估自动化安全测试能力。
\end{itemize}
& Web应用黑盒渗透测试与漏洞验证 \\
\cline{2-5}
& AgentBreaker\allowbreak~\cite{AgentBreaker}
& 2025‑07
& \begin{itemize}[tabitem]
    \item 自动测试提示注入、越狱、数据泄漏、护栏绕过、提示提取、工具滥用、数据外传和多模态注入风险；
    \item 根据已经发现的系统行为和漏洞，自适应规划后续测试攻击；
    \item 对发现结果进行严重性评分，并按照OWASP和MITRE框架进行分类。
\end{itemize}
& AI系统自动化红队测试与漏洞评估 \\
\cline{2-5}
& CAI Framework~\cite{CAI,CAI_paper}
& 2025‑06
& \begin{itemize}[tabitem]
    \item 提供面向攻防安全自动化的开源AI智能体框架；
    \item 支持构建专用安全Agent，覆盖侦察、漏洞发现、利用验证和安全评估；
    \item 结合工具集成、人类监督和guardrails，用于授权Bug Bounty、CTF 和安全测试。
\end{itemize}
& 定制化攻防演练 \\
\cline{2-5}
& Nebula~\cite{Nebula}
& 2025‑05
& \begin{itemize}[tabitem]
    \item 将AI模型集成到命令行渗透测试工作流中；
    \item 自动化记录、分类安全发现，并提供实时AI分析建议；
    \item 支持调用外部工具结果进行漏洞分析、报告整理和测试过程留痕。
\end{itemize}
& 渗透测试、漏洞分析与测试记录 \\
% \cline{2-5}
% & Pentest GPT~\cite{Deng2024PentestGPT}
% & 2024‑09
% & \begin{itemize}[tabitem]
%     \item 基于 LLM 辅助渗透测试过程的任务拆解、工具使用和结果解释。
%     \item 通过多模块交互缓解长链路渗透测试中的上下文丢失问题。
%     \item 支持真实测试目标和 CTF 场景中的攻击路径推演与验证。
% \end{itemize}
% & 自动化渗透测试、真实目标安全评估与CTF \\
\hline
\SetCell[r=4]{c} AI增强传统安全工具
& Ghidra + GhidraMCP \allowbreak~\cite{Crawford2026GhidraMCPAttack}
& 2026‑04
& \begin{itemize}[tabitem]
    \item 研究 GhidraMCP等工具增强的逆向AI智能体如何自动化二进制分析；
    \item 通过对反编译代码植入隐蔽指令，评估LLM驱动逆向流程的提示注入风险；
    \item 分析攻击如何误导disassembly/decompilation解释并污染自动化分析结果。
\end{itemize}
& 可执行二进制与恶意软件逆向分析、逆向AI Agent安全研究 \\
\cline{2-5}
& SQLMap + FastMCP \allowbreak~\cite{SQLMap,FastMCP}
& 2025‑02
& \begin{itemize}[tabitem]
    \item 将sqlmap的SQL注入检测、利用验证、数据库指纹识别等能力封装为MCP工具；
    \item 通过FastMCP为智能体提供标准化工具接口、参数校验与调用结果返回；
    \item 支持智能体在授权测试场景中编排SQL注入检测流程，并辅助整理风险结果。
\end{itemize}
& Web应用SQL注入检测与MCP工具集成 \\
% \cline{2-5}
% & Nuclei（ + AI 模板扩展） ~\cite{Nuclei}
% & 2024‑04
% & \begin{itemize}[tabitem]
%     \item 基于 YAML 模板进行高性能漏洞扫描与自定义检测。
%     \item 通过 AI prompt 辅助生成/运行检测模板。
%     \item 用于 Web、网络、DNS、云配置等多协议资产的批量扫描。
% \end{itemize}
% & 批量漏洞扫描 \\
% \cline{2-5}
% & BloodHound CE~\cite{BloodHound}
% & 2023‑11
% & \begin{itemize}[tabitem]
%     \item 基于图数据库分析身份、权限和访问关系。
%     \item 帮助红队发现复杂攻击路径，帮助蓝队识别并缓解权限风险。
%     \item 映射不同身份平台中的权限关系和潜在攻击路径
% \end{itemize}
% & 身份与权限关系分析、攻击路径发现与风险缓解 \\
\hline
\end{longtblr}
